\chapter*{\raggedright Connections:\\neighbouring differences,\\distant likenesses}
\label{connections}\addcontentsline{toc}{part}{Connections: neighbouring differences, distant likenesses}
\markboth{Connections}{Connections}
A map puts breads beside their neighbours. A table lets us rearrange them. Bring a Swedish crispbread beside a Sardinian sheet, or an Egyptian pocket beside a Venezuelan arepa, and another route through the atlas appears. We can follow the work a bread does, the surface that cooks it and the meal that gives it a reason to exist.

The following six essays make that journey through breads already met in the regions. Some connections follow documented movement of people and practice. Others begin with a likeness that becomes more interesting as its differences emerge. Shared shape alone does not establish ancestry. Each comparison leads back to the entries and their recipes, so that reading can become rolling, cooking and eating.

\begin{breadnote}{Take a question into the kitchen}
Choose two recipes, then compare one thing carefully. It might be the edge of a disc, the inside of a ring or what happens when yesterday's bread meets a sauce. Keep each recipe's ingredients and method intact for the first bake. The differences will be easier to understand when each bread has had a fair beginning.
\end{breadnote}
